% --- 1. GEOMETRIA DELLA PAGINA ---
\usepackage{geometry}
\geometry{
    a4paper,
    total={170mm,257mm},
    left=15mm,
    right=15mm,
    top=20mm,
    bottom=20mm
}

% --- 2. FONT E LEGGIBILITÀ ---
\usepackage[T1]{fontenc}
\usepackage[utf8]{inputenc}
\usepackage[italian]{babel}
\usepackage[sfdefault]{inter} 
\renewcommand{\familydefault}{\sfdefault}

% --- PACCHETTI PER IL CODICE ---
\usepackage{listings}
\usepackage{xcolor}

% Definizione colori per SQL
\definecolor{codegreen}{rgb}{0,0.6,0}
\definecolor{codegray}{rgb}{0.5,0.5,0.5}
\definecolor{codepurple}{rgb}{0.58,0,0.82}
\definecolor{backcolour}{rgb}{0.95,0.95,0.92}

% Configurazione dello stile SQL
%
\lstdefinestyle{sqlstyle}{
    backgroundcolor=\color{backcolour},   
    commentstyle=\color{codegreen},
    keywordstyle=\color{magenta},
    numberstyle=\tiny\color{codegray},
    stringstyle=\color{codepurple},
    basicstyle=\ttfamily\footnotesize,
    breakatwhitespace=false,         
    breaklines=true,                 
    captionpos=b,                    
    keepspaces=true,                 
    numbers=left,                    
    numbersep=5pt,                  
    showspaces=false,                
    showstringspaces=false,
    showtabs=false,                  
    tabsize=2,
    language=SQL
}

% Comando per impostare lo stile di default
\lstset{style=sqlstyle}

% Pacchetto per la personalizzazione delle liste
\usepackage{enumitem}

% Configurazione globale per le liste
\setlist[itemize]{topsep=0pt, itemsep=2pt, parsep=2pt}
% Personalizzazione del secondo livello (sub-lista)
\setlist[itemize,2]{label=$\circ$, leftmargin=2em}

\usepackage[most]{tcolorbox}
\tcbuselibrary{listings} % Importante per integrare listings

\newtcblisting{sqlbox}[2][]{
    arc=0mm,                          % Angoli netti per stile informatico
    enhanced,
    listing only,                     % Visualizza solo il codice
    colback=gray!5!white,             % Sfondo grigio chiaro
    colframe=magenta!75!black,        % Colore bordo (abbinato ai keyword SQL)
    fonttitle=\bfseries\ttfamily,
    title={#2},                       % Titolo nell'etichetta
    attach boxed title to top left={yshift=-2mm, xshift=2mm},
    boxed title style={colback=magenta!75!black},
    listing options={
        language=SQL,
        style=sqlstyle,               % Usa lo stile definito in precedenza
        breaklines=true
    },
    #1
}